\documentclass[../dayan_api_reference.tex]{subfiles}

\begin{document}

\section{HAL API}\label{sec:api-hal}

HAL 向 Server 和 Service 提供设备控制、物理内存管理和传输授权；向 Runtime 提供可在
用户进程内执行的数据搬运接口。公开头文件位于
\code{host/driver/hal/include/}，包括 \code{npu_hal_config.h}、\code{hal.h} 和
\code{dataplane.h}。内部能力表、具体后端的配置以及后端实现不属于公开 API。

HAL 由两个相互独立的部分组成：

\begin{itemize}
  \item 控制面运行在 Server 进程中，负责打开设备、管理物理内存、访问寄存器、处理中断，
        并为每次 Host 与 Device 之间的传输生成令牌。
  \item 数据面运行在使用 Runtime 的用户进程中。数据提供器根据令牌直接访问设备文件，
        不要求 Server 读取或暂存用户数据。
\end{itemize}

后端是控制面的平台实现，数据提供器是与该后端配套的用户进程传输实现。两者通过数值
标识和不透明字段配对。通用 HAL 不解释后端写入的不透明字段。

\subsection{公开类型与常量}

\begin{dayantable}{HAL 核心类型}{tab:api-hal-types}{X[34,l] X[66,l]}
  类型 & 用途 \\
  \code{npu_hal_t} & 已打开的 HAL 实例，保存所选后端、设备拓扑、内存管理状态和首个不可恢复的 I/O 错误。 \\
  \code{npu_hal_allocation_t} & 一段由 HAL 管理的设备物理内存。该类型是不透明句柄。 \\
  \code{npu_hal_core_t} & 用 Cluster 编号和 Core 编号表示一个物理执行位置。 \\
  \code{npu_hal_memory_class_t} & 内存类别：普通内存、一致性内存或 Core 私有 URAM。 \\
  \code{npu_hal_memory_scope_t} & 内存可见范围：整个设备、一个 Cluster 或一个 Core。 \\
  \code{npu_hal_memory_usage_t} & 内存用途，例如用户 Buffer、固件、算子、堆、启动描述符、共享区、URAM、性能采样区或 Atomic Cache 后备区。 \\
  \code{npu_hal_memory_region_t} & 一段可分配的物理地址区间，以及它的类别、可见范围和归属位置。 \\
  \code{npu_hal_backend_options_t} & 所有后端配置共有的头部，包含配置版本和后端标识。 \\
  \code{npu_hal_options_t} & HAL 通用配置，包含设备拓扑、普通内存默认容量、内存区域表和后端配置指针。 \\
  \code{npu_hal_info_t} & 已打开 HAL 的版本、拓扑、数据提供器和传输能力。 \\
  \code{npu_hal_memory_request_t} & 一次内存申请的类别、范围、用途、所有者、位置、大小、对齐和可选固定地址。 \\
  \code{npu_hal_allocation_info_t} & 已分配内存的地址和属性快照。 \\
\end{dayantable}

\begin{dayantable}{HAL 数据面类型}{tab:api-hal-data-types}{X[36,l] X[64,l]}
  类型 & 用途 \\
  \code{npu_hal_data_plane_t} & Runtime 导入设备句柄后创建的数据面实例。该类型是不透明句柄。 \\
  \code{npu_hal_data_direction_t} & 传输方向：Host 到 Device，或 Device 到 Host。 \\
  \code{npu_hal_data_export_t} & 数据提供器标识、设备句柄数量，以及仅由该提供器解释的初始化数据。 \\
  \code{npu_hal_data_token_t} & 一次传输的授权记录，包含提供器、方向、长度和后端私有定位信息。 \\
\end{dayantable}

\begin{dayantable}{HAL 枚举值与上限}{tab:api-hal-enums}{X[47,l] X[53,l]}
  API & 含义 \\
  \code{NPU_HAL_MAX_CLUSTERS} / \code{NPU_HAL_MAX_CORES_PER_CLUSTER} & 最多表示 32 个 Cluster，每个 Cluster 最多 32 个 Core。 \\
  \code{NPU_HAL_MEMORY_NORMAL} / \code{NPU_HAL_MEMORY_COHERENT} / \code{NPU_HAL_MEMORY_URAM} & 三种内存类别。 \\
  \code{NPU_HAL_MEMORY_DEVICE} / \code{NPU_HAL_MEMORY_CLUSTER} / \code{NPU_HAL_MEMORY_CORE} & 三种内存可见范围。 \\
  \code{NPU_HAL_MEMORY_BUFFER} / \code{NPU_HAL_MEMORY_FIRMWARE} / \code{NPU_HAL_MEMORY_KERNEL} & 用户 Buffer、固件和算子代码用途。 \\
  \code{NPU_HAL_MEMORY_HEAP} / \code{NPU_HAL_MEMORY_LAUNCH} / \code{NPU_HAL_MEMORY_SHARED} & Device 堆、启动描述符和共享区用途。 \\
  \code{NPU_HAL_MEMORY_KERNEL_URAM} / \code{NPU_HAL_MEMORY_PROFILE} / \code{NPU_HAL_MEMORY_ATOMIC_CACHE_BACKING} & 算子 URAM、性能采样区和 Atomic Cache 后备区用途。 \\
  \code{NPU_HAL_DATA_HOST_TO_DEVICE} / \code{NPU_HAL_DATA_DEVICE_TO_HOST} & 两种用户进程直传方向。 \\
  \code{NPU_HAL_DATA_MAX_HANDLES} & 一个数据面实例最多导入 64 个设备句柄。 \\
  \code{NPU_HAL_BACKEND_XDMA} / \code{NPU_HAL_DATA_PROVIDER_XDMA} & XDMA 控制后端和与其配套的数据提供器标识。 \\
\end{dayantable}

结构版本由 \code{NPU_HAL_ABI_VERSION}、\code{NPU_HAL_OPTIONS_VERSION}、
\code{NPU_HAL_MEMORY_REQUEST_VERSION}、\code{NPU_HAL_DATA_EXPORT_VERSION} 和
\code{NPU_HAL_DATA_TOKEN_VERSION} 标识。\code{NPU_HAL_ANY_CLUSTER} 和
\code{NPU_HAL_ANY_CORE} 表示任意物理位置；\code{NPU_HAL_TIMEOUT_INFINITE} 表示无限等待。

\subsection{生命周期与设备拓扑}

\begin{dayantable}{HAL 生命周期与查询}{tab:api-hal-life}{X[38,l] X[62,l]}
  API & 行为 \\
  \code{npu_hal_options_init} & 写入配置版本和一个 Cluster、每个 Cluster 四个 Core 的默认拓扑，不写入任何后端私有配置。 \\
  \code{npu_hal_open} & 校验通用配置，建立内存管理器并打开所选后端。未指定后端配置时使用注册表中的默认后端。 \\
  \code{npu_hal_close} & 关闭后端并释放全部内存管理和同步状态。可传入空指针。 \\
  \code{npu_hal_device_status} & 返回 0，或此前记录的首个不可恢复设备 I/O 错误。 \\
  \code{npu_hal_get_info} & 返回 HAL 版本、设备拓扑、数据提供器和传输限制。 \\
  \code{npu_hal_get_memory_region} & 按索引读取一项物理内存区域配置。 \\
  \code{npu_hal_core_valid} & 判断给定 Cluster 和 Core 编号是否存在。 \\
  \code{npu_hal_memory_capacity} & 统计指定 Core 可见的某类内存总容量。 \\
\end{dayantable}

\subsection{设备地址、寄存器、中断与启动}

原始设备地址接口只供 Service 的固件和 RPC 实现使用。Runtime 只操作逻辑 Buffer，
不会取得设备物理地址。接口中的 Cluster 参数选择一套物理地址空间，Core 参数进一步
选择该 Cluster 中的执行位置。

\begin{dayantable}{HAL 控制面访问}{tab:api-hal-control}{X[39,l] X[61,l]}
  API & 行为 \\
  \code{npu_hal_copy_to_address} & 把 Host 内存中的数据写到指定 Cluster 的设备物理地址。 \\
  \code{npu_hal_copy_from_address} & 从指定 Cluster 的设备物理地址读到 Host 内存。 \\
  \code{npu_hal_copy_address} & 借助固定大小的 Host 临时缓冲区，在同一 Cluster 内复制两段设备内存。 \\
  \code{npu_hal_register_read} & 从后端控制窗口读取寄存器。 \\
  \code{npu_hal_register_write} & 向后端控制窗口写寄存器。 \\
  \code{npu_hal_irq_status} & 读取一个 Cluster 尚未处理的 Host 中断位图。 \\
  \code{npu_hal_irq_wait} & 等待任一 Cluster 产生中断；超时返回 \code{-ETIMEDOUT} 并清空输出参数。 \\
  \code{npu_hal_irq_ack} & 清除目标 Cluster 中指定的中断位。 \\
  \code{npu_hal_irq_mailbox} & 读取指定 Core 的 Host 中断邮箱命令。 \\
  \code{npu_hal_command_submit} & 写入 Device 邮箱并触发目标 Core 的软件中断。 \\
  \code{npu_hal_wait_command_consumed} & 等待目标 Core 取走邮箱命令。超时时间以毫秒为单位。 \\
  \code{npu_hal_boot_core} & 等待旧命令结束，提交 BL1 入口，再等待 BL0 取走启动命令。 \\
\end{dayantable}

设备文件返回 \code{EIO}、\code{ENODEV}、\code{EBADF} 等不可恢复错误时，HAL 会保存
首个错误。后续控制操作直接返回该错误，Server 据此停止服务。

\subsection{物理内存管理}

\begin{dayantable}{HAL 内存管理 API}{tab:api-hal-memory}{X[39,l] X[61,l]}
  API & 行为 \\
  \code{npu_hal_memory_request_init} & 初始化一个设备级普通 Buffer 申请，默认按 64 字节对齐。 \\
  \code{npu_hal_memory_alloc} & 在类别、可见范围和位置均匹配的区域中查找一段空闲地址。 \\
  \code{npu_hal_memory_reserve} & 在匹配的区域中占用调用方指定的固定地址范围。 \\
  \code{npu_hal_memory_free} & 释放一段 HAL 内存分配；调用后原句柄失效。 \\
  \code{npu_hal_memory_info} & 读取一段已分配内存的地址和属性。 \\
  \code{npu_hal_memory_visible} & 判断一段已分配内存是否对指定 Core 可见。 \\
  \code{npu_hal_address_valid} & 判断一个原始地址范围是否落在指定 Core 可见的区域内。 \\
  \code{npu_hal_memory_address} & 取得一段已分配内存对指定 Core 可见的设备地址。 \\
  \code{npu_hal_memory_write} & 校验句柄、访问位置和范围后，把 Host 数据写入已分配内存。 \\
  \code{npu_hal_memory_read} & 校验句柄、访问位置和范围后，把已分配内存读到 Host。 \\
\end{dayantable}

\subsection{令牌化数据面}

数据面的目标是让用户进程直接访问设备，避免把大块数据先送到 Server 再转发。当前
XDMA 数据提供器不创建、不连接也不读写 Socket。它运行在 Runtime 所在的用户进程中，
直接对 \code{/dev} 或 \code{/tmp/virtual_npu} 下的 H2C/C2H 设备文件执行
\code{pread}/\code{pwrite}。

Unix 域 Socket 只属于 Runtime 与 Server 之间的通用协议。建立会话时，Server 使用
\code{SCM_RIGHTS} 把已经打开的设备文件描述符交给用户进程。\code{SCM_RIGHTS} 传递的
是文件描述符，不是待搬运的数据。之后的授权请求和完成通知仍是短控制消息；Host Buffer
中的数据不会进入 Socket，也不会进入 Server 缓冲区。因此这里减少的是中转拷贝，
并不承诺内核和设备驱动层面的零拷贝。

\begin{dayantable}{HAL 数据面 API}{tab:api-hal-data}{X[42,l] X[58,l]}
  API & 行为 \\
  \code{npu_hal_data_plane_export} & Server 侧要求后端填写数据提供器标识，并复制用户进程直接访问设备所需的文件描述符。该函数只生成导出结果，不负责 Socket 发送。 \\
  \code{npu_hal_transfer_prepare} & Server 侧校验内存归属、可见范围、偏移、长度和方向，再为这一笔传输生成令牌。 \\
  \code{npu_hal_transfer_finish} & Server 侧核对用户进程返回的完整令牌，把数据提供器的传输结果交给后端完成收尾，并返回处理结果。 \\
  \code{npu_hal_data_plane_import} & Runtime 侧根据提供器标识选择实现，接管导入的设备文件描述符，并创建数据面实例。 \\
  \code{npu_hal_data_plane_close} & 关闭数据面实例持有的设备文件描述符并释放内部状态。 \\
  \code{npu_hal_data_plane_transfer} & Runtime 侧把令牌和用户缓冲区交给数据提供器。XDMA 提供器据此选择 H2C 或 C2H 文件、Cluster 和设备偏移，直接执行读写。 \\
\end{dayantable}

一次 H2D 或 D2H 传输按以下顺序执行：

\begin{enumerate}
  \item Runtime 只把逻辑 Buffer、偏移和长度提交给 Server，不发送用户数据。
  \item Server 检查对象所有权和物理内存范围。命令可以执行时，HAL 生成一次性令牌。
  \item Runtime 的传输线程取得令牌，调用匹配的数据提供器。XDMA 提供器直接读写设备文件。
  \item Runtime 把原令牌和传输结果交还 Server。Server 核对后完成 Event，并允许相关内存被释放或复用。
\end{enumerate}

导出描述和令牌中的后端私有字段只由匹配的数据提供器解释。通用 Runtime 不推导物理
地址，也不需要知道 Cluster 到设备文件的映射。其它 HAL 后端应遵循同一原则：控制面
签发令牌，用户进程中的数据提供器使用导入的设备资源直接传输；具体实现可以使用
\code{read/write}、\code{ioctl} 或 \code{mmap}，但不应让 Server 代理数据本身。

令牌是 Runtime、Server 和 HAL 之间的协议授权与生命周期凭据。当前 XDMA 驱动把完整
设备文件描述符交给用户进程，因此令牌本身不是针对恶意进程的内核级访问隔离；进程权限
和设备文件权限仍是安全边界。

\end{document}
